\documentclass[pdflatex,sn-mathphys-num]{sn-jnl}

\usepackage{graphicx}
\usepackage{multirow}
\usepackage{amsmath,amssymb,amsfonts}
\usepackage{amsthm}
\usepackage{mathrsfs}
\usepackage[title]{appendix}
\usepackage{xcolor}
\usepackage{textcomp}
\usepackage{manyfoot}
\usepackage{booktabs}
\usepackage{algorithm}
\usepackage{algorithmicx}
\usepackage{algpseudocode}
\usepackage{listings}

% Instrucciones para el estudiante: encabezado en azul + texto en cursiva,
% sin caja, para no dejar espacio en blanco extra en la página.
\newenvironment{instruccion}{%
  \par\noindent\textbf{\color{blue!55!black}INSTRUCCIONES PARA ESTA SECCIÓN}%
  \par\vspace{2pt}\noindent\itshape\small
}{%
  \par\normalfont\vspace{6pt}
}

% Ejemplo completamente desarrollado: encabezado en verde + texto normal,
% sin caja, para no dejar espacio en blanco extra en la página.
\newenvironment{ejemplo}{%
  \par\noindent\textbf{\color{green!35!black}EJEMPLO DESARROLLADO (así debe verse tu sección terminada)}%
  \par\vspace{2pt}\noindent\small
}{%
  \par\normalfont\vspace{6pt}
}

\raggedbottom

\begin{document}

\title[Ejemplo: CNN Podada para Retinopatía Diabética]{Un Ensamble Ligero Basado en Poda Dinámica de Filtros para la Clasificación de Retinopatía Diabética: Reduciendo Error y Tiempo de Cómputo frente a Modelos SOTA}

\author*[1]{\fnm{\textcolor{blue}{Lihki Rubio}}}\email{lihkir@uninorte.edu.co}

\author[1]{\fnm{Carlos} \sur{Rodríguez}}\email{carlos.rodriguez@universidad.edu}

\author[1]{\fnm{Laura} \sur{Gómez}}\email{laura.gomez@universidad.edu}

\affil*[1]{\orgdiv{Departamento de Matemáticas, Física y Ciencia de Datos}, \orgname{Universidad del Norte}, \orgaddress{\city{Barranquilla}, \country{Colombia}}}

\abstract{La retinopatía diabética (RD) es una de las principales causas de ceguera prevenible en adultos, y su diagnóstico automatizado mediante imágenes de fondo de ojo suele apoyarse en redes convolucionales profundas de alto costo computacional. En este trabajo se propone una CNN ligera con poda dinámica de filtros durante el entrenamiento, evaluada sobre el conjunto de datos público APTOS 2019 (3{,}662 imágenes, 5 clases de severidad). El modelo propuesto se comparó contra dos referencias SOTA vistas en clase, ResNet-50 y EfficientNet-B0, bajo el mismo esquema de validación cruzada de 5 particiones. Los resultados muestran que el modelo propuesto alcanza un F1-score macro de 0.912 $\pm$ 0.008, superando a EfficientNet-B0 (0.887 $\pm$ 0.011) y a ResNet-50 (0.869 $\pm$ 0.014), mientras reduce el tiempo de entrenamiento en un 34\% y el tiempo de inferencia por imagen en un 41\% respecto al mejor baseline. Estos resultados sugieren que la poda dinámica de filtros permite mantener, e incluso mejorar, la capacidad discriminativa del modelo mientras se reduce significativamente el costo computacional, lo cual es relevante para su despliegue en entornos clínicos con recursos limitados.}

\keywords{Clasificación de imágenes médicas, Redes neuronales convolucionales, Poda de filtros, Retinopatía diabética, F1-score}

\maketitle

\begin{instruccion}
El título debe indicar (a) el tipo de problema, clasificación o regresión, (b) el dominio o dataset, y (c) qué hace novedoso a tu modelo. El abstract va en un único párrafo de 150 a 250 palabras, sin citas ni ecuaciones, y debe responder en orden: problema y dataset, idea del método propuesto, resultados cuantitativos frente a los modelos SOTA (métrica y tiempo de cómputo), y una conclusión breve. Las palabras clave (4 a 6) deben incluir el nombre de tu técnica, el tipo de tarea, el dominio y la métrica principal. El bloque de abajo es un ejemplo completo y realista de cómo debe verse tu título, autores y abstract terminados; reemplázalo por los tuyos manteniendo esta misma estructura y nivel de detalle.
\end{instruccion}

\section{Introducción}\label{sec:intro}

\begin{instruccion}
Escribe entre 1 y 1.5 páginas que respondan: por qué es importante el problema, qué limitaciones tienen los modelos SOTA actuales que motivan tu propuesta, y cuál es tu contribución principal. Cierra con una oración que resuma tu hipótesis, por ejemplo: ``En este trabajo proponemos... y demostramos que supera a los modelos SOTA X, Y en la métrica Z, reduciendo además el tiempo de cómputo en W\%.'' Cita aquí los trabajos de los modelos SOTA usando \texttt{\textbackslash cite\{\}}, con las entradas correspondientes en tu archivo \texttt{sn-bibliography.bib}.
\end{instruccion}

\begin{ejemplo}
El cribado automatizado de retinopatía diabética a partir de imágenes de fondo de ojo ha mostrado resultados prometedores con arquitecturas profundas como ResNet-50 \cite{bib1} y EfficientNet-B0 \cite{bib2}. Sin embargo, estos modelos requieren un alto costo computacional en entrenamiento e inferencia, lo cual limita su despliegue en centros de salud con hardware restringido. Trabajos recientes en poda de redes neuronales \cite{bib3} sugieren que es posible eliminar filtros redundantes sin afectar de forma significativa la exactitud del modelo. Partiendo de esta idea, en este trabajo proponemos una CNN con poda dinámica de filtros guiada por la magnitud de los gradientes durante el entrenamiento, y demostramos que supera a ResNet-50 y EfficientNet-B0 en F1-score macro, reduciendo simultáneamente el tiempo de entrenamiento y de inferencia.
\end{ejemplo}

\section{Metodología}\label{sec:metodologia}

\begin{instruccion}
Esta es la sección más importante del proyecto: debe permitir que cualquier persona reproduzca tu experimento leyendo únicamente el texto. Organízala en las cuatro subsecciones siguientes y no omitas ninguna; si algo no aplica, indícalo y explica por qué.
\end{instruccion}

\subsection{Descripción del problema y del conjunto de datos}\label{subsec:datos}

\begin{instruccion}
Esta subsección debe contener tres partes, en este orden: \textbf{(1) descripción general del dataset}, \textbf{(2) proceso ETL} y \textbf{(3) Análisis Exploratorio de Datos (EDA) exhaustivo}. No omitas ninguna.

\textbf{(1) Descripción general.} Indica: nombre y fuente del dataset (con enlace o cita), número de instancias y de variables, tipo de variable objetivo, y balance de clases si es clasificación.

\textbf{(2) ETL (Extract, Transform, Load).} Describe explícitamente las tres etapas:
\begin{itemize}
\item \textit{Extract:} de dónde y en qué formato se obtuvieron los datos crudos (CSV, imágenes, API, base de datos, etc.).
\item \textit{Transform:} limpieza de datos (duplicados, valores faltantes y cómo se trataron: eliminación, imputación por media/mediana/moda, imputación por modelo, etc.), corrección de tipos de datos, tratamiento de outliers, codificación de variables categóricas (one-hot, ordinal, etc.), normalización/estandarización, y técnicas de aumento de datos si aplica.
\item \textit{Load:} cómo y en qué formato queda el dataset final listo para el modelado (ej. tensores, dataframe procesado, tamaño final tras la limpieza), y la partición usada (por ejemplo 70/15/15 o validación cruzada k-fold indicando k), especificando también la semilla aleatoria usada para garantizar reproducibilidad.
\end{itemize}

\textbf{(3) EDA exhaustivo.} El Análisis Exploratorio de Datos debe ser exhaustivo, no un simple vistazo. Como mínimo debe incluir:
\begin{itemize}
\item Estadística descriptiva de las variables relevantes (media, mediana, desviación estándar, mínimos, máximos, valores faltantes por variable).
\item Análisis univariado: distribución de cada variable relevante y de la variable objetivo (histogramas, gráficos de barras, o su equivalente para imágenes/texto, como distribución de brillo/tamaño o longitud de secuencias).
\item Análisis bivariado o multivariado: relación entre las variables predictoras y la variable objetivo, y matriz de correlación entre variables si aplica.
\item Detección visual de valores atípicos (outliers), por ejemplo con boxplots.
\item Visualización del balance de clases (si es clasificación) o de la distribución de la variable objetivo (si es regresión).
\item Al menos una figura (usando el entorno \texttt{figure}) que resuma un hallazgo relevante del EDA.
\item Un párrafo de conclusiones del EDA que explique qué decisiones de preprocesamiento (de la etapa Transform del ETL) se tomaron \textbf{como consecuencia directa} de lo observado en el EDA (por ejemplo, aplicar aumento de datos porque el EDA reveló desbalance de clases).
\end{itemize}
\end{instruccion}

\begin{ejemplo}
\textbf{Descripción general.} Se utilizó el conjunto público APTOS 2019 \cite{bib4}, compuesto por 3{,}662 imágenes de fondo de ojo etiquetadas en cinco niveles de severidad de retinopatía diabética (0: sin RD, hasta 4: RD proliferativa).

\textbf{ETL.} \textit{Extract:} las imágenes y sus etiquetas se descargaron en formato \texttt{.png} junto con un archivo \texttt{.csv} de metadatos desde el repositorio oficial de la competencia \cite{bib4}. \textit{Transform:} se eliminaron 12 imágenes duplicadas y 5 imágenes corruptas detectadas mediante verificación de integridad del archivo; las imágenes se redimensionaron a $224\times224$ píxeles y se normalizaron con la media y desviación estándar de ImageNet; no se registraron valores faltantes en las etiquetas. \textit{Load:} el dataset final, con 3{,}645 imágenes tras la limpieza, se almacenó como tensores listos para entrenamiento, particionado con validación cruzada estratificada de 5 particiones (5-fold) y un conjunto de prueba independiente (15\% de los datos, no usado durante el ajuste de hiperparámetros), usando semilla aleatoria fija (\texttt{seed=42}) para garantizar reproducibilidad.

\textbf{EDA exhaustivo.} El análisis univariado de la variable objetivo reveló un fuerte desbalance de clases: la clase 0 (sin RD) representa el 49\% de las imágenes, mientras que la clase 4 (RD proliferativa) representa solo el 5\%, como se muestra en la Figura~\ref{fig:eda}. El análisis del brillo promedio por imagen mostró una distribución aproximadamente normal, sin outliers extremos que sugirieran errores de captura. El análisis bivariado (brillo promedio vs. severidad) no mostró una correlación clara, por lo que el brillo no se usó como variable auxiliar. Como consecuencia directa de estos hallazgos, se decidió aplicar aumento de datos (rotación $\pm$15°, volteo horizontal y variación de brillo) únicamente sobre las clases minoritarias (2, 3 y 4), en lugar de aplicarlo de forma uniforme a todo el dataset.
\end{ejemplo}

% NOTA TÉCNICA: recuerda que "figure" es un entorno flotante y, al igual que
% "algorithm" y "table", conviene mantenerlo fuera de los entornos
% "instruccion"/"ejemplo" (que aquí solo aplican formato de párrafo, no cajas),
% para evitar problemas de flotación y mantener una diagramación limpia.
% Por eso la figura del EDA va aquí, fuera del entorno "ejemplo".
\begin{figure}[h]
\centering
% \includegraphics[width=0.65\textwidth]{distribucion_clases_eda.png}
\caption{Distribución de clases del dataset APTOS 2019 obtenida durante el EDA, evidenciando el desbalance hacia la clase 0 (sin RD).}
\label{fig:eda}
\end{figure}

\subsection{Modelos SOTA de referencia (baselines)}\label{subsec:sota}

\begin{instruccion}
Debes comparar tu modelo propuesto contra \textbf{todos} los modelos vistos en la clase de Machine Learning, incluyendo tanto los modelos clásicos como los modelos más avanzados/profundos que se hayan estudiado en el curso. Para cada modelo de referencia: cítalo, explica brevemente por qué se considera un buen punto de comparación para este tipo de problema, \textbf{define matemáticamente el modelo de forma resumida} (su formulación o función objetivo principal, usando el entorno \texttt{equation}), agregando la(s) referencia(s) bibliográficas relevantes de dicha formulación, e indica los hiperparámetros usados (idealmente ajustados con Grid Search o Random Search, señalando el rango explorado). El objetivo del proyecto es que tu modelo novedoso propuesto sea superior a \textbf{todos} estos modelos simultáneamente en (a) la métrica de error o precisión, y (b) el tiempo de cómputo (entrenamiento y/o inferencia). Si tu modelo no logra superar a alguno de ellos en alguna de las dos dimensiones, debes reportarlo igualmente de forma honesta y discutirlo en Resultados y Conclusiones.
\end{instruccion}

\begin{ejemplo}
Se compararon todos los modelos vistos en el curso de Machine Learning: \textbf{Regresión Logística} \cite{bib5}, como línea base lineal simple; \textbf{Máquinas de Vectores de Soporte (SVM)} con kernel RBF \cite{bib6}; \textbf{Random Forest} \cite{bib7}, por su robustez ante datos desbalanceados; \textbf{XGBoost} \cite{bib8}, como referencia de boosting de árboles ampliamente usada en competencias; \textbf{ResNet-50} \cite{bib1}, preentrenada en ImageNet y afinada (fine-tuning) sobre APTOS 2019, estándar de facto en clasificación de imágenes médicas; y \textbf{EfficientNet-B0} \cite{bib2}, también preentrenada, incluida por su balance entre exactitud y número de parámetros. Para los modelos clásicos (Regresión Logística, SVM, Random Forest, XGBoost) se usaron como entrada características extraídas con un extractor preentrenado congelado, y se ajustaron sus hiperparámetros principales mediante búsqueda en rejilla (por ejemplo, $C \in \{0.1, 1, 10\}$ para SVM, profundidad $\in \{4, 6, 8\}$ y número de árboles $\in \{100, 300, 500\}$ para Random Forest y XGBoost). Para ResNet-50 y EfficientNet-B0 se realizó una búsqueda en rejilla sobre tasa de aprendizaje ($\{10^{-3}, 10^{-4}, 10^{-5}\}$) y tamaño de lote ($\{16, 32, 64\}$), seleccionando en todos los casos la combinación con mejor F1-score macro en validación; los modelos basados en redes neuronales se entrenaron durante 50 épocas con early stopping (paciencia de 8 épocas).

A modo de ejemplo, a continuación se presenta la definición matemática resumida de uno de los modelos de referencia, la \textbf{SVM con kernel RBF} \cite{bib6}. Dado un conjunto de entrenamiento $\{(x_i, y_i)\}_{i=1}^{n}$ con $y_i \in \{-1, +1\}$, la SVM busca el hiperplano de separación óptimo resolviendo el siguiente problema de optimización con margen suave:

\begin{equation}
\min_{w, b, \xi} \; \frac{1}{2}\lVert w \rVert^2 + C\sum_{i=1}^{n}\xi_i \quad \text{sujeto a} \quad y_i(w^\top \phi(x_i) + b) \geq 1 - \xi_i, \;\; \xi_i \geq 0,
\label{eq:svm_primal}
\end{equation}

\noindent donde $\phi(\cdot)$ es el mapeo implícito al espacio de características inducido por el kernel, $\xi_i$ son variables de holgura y $C$ controla el trade-off entre margen y error de clasificación. Resolviendo el problema dual y aplicando el truco del kernel, la función de decisión resultante es

\begin{equation}
f(x) = \operatorname{sign}\left(\sum_{i=1}^{n} \alpha_i y_i \, k(x_i, x) + b\right), \qquad k(x_i, x) = \exp\left(-\gamma \lVert x_i - x \rVert^2\right),
\label{eq:svm_dual}
\end{equation}

\noindent donde $\alpha_i$ son los multiplicadores de Lagrange obtenidos al resolver el problema dual y $k(x_i, x)$ es el kernel RBF (Gaussiano), con $\gamma$ como hiperparámetro que controla la influencia de cada punto de entrenamiento \cite{bib6}. Cada uno de los demás modelos de referencia (Regresión Logística, Random Forest, XGBoost, ResNet-50, EfficientNet-B0) debe presentarse con el mismo nivel de detalle: su formulación matemática resumida y la(s) referencia(s) correspondiente(s).
\end{ejemplo}

\subsection{Modelo propuesto}\label{subsec:propuesto}

\begin{instruccion}
Explica tu modelo novedoso con el mayor detalle técnico posible: arquitectura o algoritmo, ecuaciones relevantes, y pseudocódigo. Incluye la justificación teórica de por qué esperas que esta idea mejore la precisión o reduzca el error, y/o reduzca el tiempo de cómputo frente a los modelos de la subsección anterior. No basta con decir que funcionó: debe existir una razón lógica o teórica detrás del diseño.
\end{instruccion}

\begin{ejemplo}
Se propone una CNN ligera basada en una arquitectura tipo ResNet reducida (18 capas), a la que se aplica poda dinámica de filtros durante el entrenamiento: en cada época, los filtros cuya norma $L_1$ de gradiente acumulado esté por debajo del percentil 10 se marcan como candidatos a eliminación, y se eliminan definitivamente si permanecen por debajo de ese umbral durante 3 épocas consecutivas. Formalmente, para el filtro $k$ de la capa $l$, la importancia se define como

\begin{equation}
I^{(l)}_k = \frac{1}{E}\sum_{e=1}^{E}\left\lVert \nabla_{W^{(l)}_k} \mathcal{L}_e \right\rVert_1,
\label{eq:importancia}
\end{equation}

\noindent donde $E$ es la ventana de épocas evaluada y $\mathcal{L}_e$ la función de pérdida en la época $e$. La hipótesis es que, al eliminar progresivamente los filtros con menor contribución al gradiente, se reduce la redundancia de la red sin afectar su capacidad discriminativa, lo que en clasificación desbalanceada como RD puede incluso reducir el sobreajuste hacia la clase mayoritaria y mejorar el F1-score macro, además de reducir el número de operaciones de punto flotante (FLOPs) y, por tanto, el tiempo de entrenamiento e inferencia. El procedimiento completo se resume en el Algoritmo~\ref{alg:propuesto}, mostrado a continuación.
\end{ejemplo}

% NOTA TÉCNICA: los entornos "flotantes" (algorithm, figure, table) van fuera
% de "instruccion"/"ejemplo", igual que antes, para mantener una diagramación
% limpia y evitar problemas de flotación en LaTeX.
\begin{algorithm}
\caption{Entrenamiento con poda dinámica de filtros}\label{alg:propuesto}
\begin{algorithmic}[1]
\Require Datos de entrenamiento $D$, red inicial $f_\theta$, ventana $E$, umbral $\tau$ (percentil 10)
\Ensure Modelo entrenado y podado $f_{\theta'}$
\For{cada época $e = 1, \dots, N$}
    \State Entrenar $f_\theta$ una época sobre $D$ y registrar gradientes por filtro
    \If{$e \bmod E = 0$}
        \State Calcular $I^{(l)}_k$ con la Ecuación~\eqref{eq:importancia} para cada filtro
        \State Marcar como candidatos los filtros con $I^{(l)}_k$ bajo el percentil $\tau$
        \State Eliminar los filtros marcados en 3 evaluaciones consecutivas
    \EndIf
\EndFor
\State \Return $f_{\theta'}$
\end{algorithmic}
\end{algorithm}

\subsection{Configuración experimental}\label{subsec:config}

\begin{instruccion}
Especifica: métricas de evaluación y por qué son adecuadas para tu problema, hardware usado, software y versiones, y cómo mediste el tiempo de cómputo (entrenamiento y/o inferencia, en segundos, y en qué hardware). Indica si repetiste cada experimento varias veces con distintas semillas para reportar promedio y desviación estándar en Resultados; esto es muy recomendable para poder afirmar honestamente que tu modelo es mejor.
\end{instruccion}

\begin{ejemplo}
Dado el desbalance de clases, se usaron F1-score macro y exactitud balanceada como métricas principales, complementadas con la matriz de confusión. Todos los experimentos se ejecutaron en una GPU NVIDIA T4 (16GB), usando Python 3.10, PyTorch 2.1 y scikit-learn 1.3. Cada modelo se entrenó 10 veces con semillas distintas ($\{1, \dots, 10\}$) para reportar promedio $\pm$ desviación estándar. El tiempo de entrenamiento se midió como el tiempo total hasta el criterio de early stopping, y el tiempo de inferencia como el promedio por imagen sobre el conjunto de prueba (1{,}000 imágenes), ambos con la misma GPU para asegurar una comparación justa.
\end{ejemplo}

\section{Resultados}\label{sec:resultados}

\begin{instruccion}
El modelo propuesto \textbf{debe superar a todos los modelos SOTA de referencia simultáneamente en (a) la métrica de precisión/error y (b) el tiempo de cómputo} (entrenamiento y/o inferencia). Este es un requisito del proyecto, no un resultado opcional: si tu modelo no logra superar a todos los modelos SOTA en ambas dimensiones, \textbf{no se considera un resultado válido para la entrega}; en ese caso debes volver a la Metodología, replantear o modificar tu idea, y proponer otro modelo original hasta lograrlo, en lugar de reportar y discutir un resultado que no cumple el objetivo. Una vez que tengas resultados que sí cumplan esta condición, presenta lo siguiente: una tabla comparativa con tu modelo y todos los SOTA de referencia mostrando la(s) métrica(s) principal(es) y el tiempo de cómputo, con promedio $\pm$ desviación estándar; al menos una figura de apoyo (curva de aprendizaje, matriz de confusión, curva ROC, etc.); y un breve análisis de si la diferencia observada frente a cada modelo SOTA es estadísticamente significativa, idealmente con una prueba estadística como t-test pareado o Wilcoxon. En el texto no repitas los números exactos de la tabla, resalta lo más relevante.
\end{instruccion}

\begin{ejemplo}
La Tabla~\ref{tab:resultados} resume el desempeño de todos los modelos comparados. El modelo propuesto obtuvo el F1-score macro más alto (0.912 $\pm$ 0.008) de todo el conjunto de modelos evaluados, superando al mejor modelo clásico (XGBoost, 0.851) por 6.1 puntos porcentuales y al mejor modelo profundo (EfficientNet-B0, 0.887) por 2.5 puntos, mientras que su tiempo de cómputo (entrenamiento e inferencia) fue el más bajo de todos los modelos basados en redes neuronales y comparable al de los modelos clásicos más livianos. Un test t pareado sobre las 10 corridas confirmó que la diferencia en F1-score frente a cada uno de los demás modelos es estadísticamente significativa ($p < 0.05$ en todos los casos). La Figura~\ref{fig:resultados} muestra la matriz de confusión del modelo propuesto sobre el conjunto de prueba, donde se observa que la mayor confusión ocurre entre las clases 1 y 2 (RD leve y moderada), un patrón consistente con lo reportado en la literatura para este dataset.
\end{ejemplo}

\begin{table}[h]
\caption{Comparación del modelo propuesto contra todos los modelos vistos en el curso de Machine Learning. Promedio $\pm$ desviación estándar de 10 corridas sobre el conjunto de prueba.}\label{tab:resultados}
\begin{tabular}{@{}lccc@{}}
\toprule
Modelo & F1-score macro & Exactitud balanceada & Tiempo entren. (min) / infer. (ms) \\
\midrule
Regresión Logística & 0.712 $\pm$ 0.019 & 0.705 $\pm$ 0.017 & 1.2 / 1.1 \\
SVM (kernel RBF)     & 0.764 $\pm$ 0.016 & 0.758 $\pm$ 0.015 & 8.4 / 4.3 \\
Random Forest        & 0.823 $\pm$ 0.013 & 0.815 $\pm$ 0.012 & 6.7 / 3.8 \\
XGBoost               & 0.851 $\pm$ 0.012 & 0.844 $\pm$ 0.011 & 9.9 / 5.2 \\
ResNet-50             & 0.869 $\pm$ 0.014 & 0.861 $\pm$ 0.012 & 42.1 / 18.4 \\
EfficientNet-B0       & 0.887 $\pm$ 0.011 & 0.879 $\pm$ 0.010 & 35.6 / 12.9 \\
\textbf{Modelo propuesto} & \textbf{0.912 $\pm$ 0.008} & \textbf{0.905 $\pm$ 0.009} & \textbf{23.5 / 7.6} \\
\botrule
\end{tabular}
\end{table}

\begin{figure}[h]
\centering
% \includegraphics[width=0.65\textwidth]{matriz_confusion.png}
\caption{Matriz de confusión del modelo propuesto sobre el conjunto de prueba de APTOS 2019. Las clases 1 y 2 concentran la mayor confusión, en línea con la dificultad conocida de distinguir grados leves de RD.}
\label{fig:resultados}
\end{figure}

\section{Conclusiones}\label{sec:conclusiones}

\begin{instruccion}
Escribe entre 1 y 2 párrafos que respondan en este orden: si se cumplió el objetivo de superar a los modelos SOTA en precisión/error y/o en tiempo de cómputo (sé honesto si solo mejoraste una de las dos cosas); cuáles son las limitaciones de tu estudio; y qué trabajo futuro propondrías. No repitas los números exactos de la tabla de resultados, interpreta su significado.
\end{instruccion}

\begin{ejemplo}
El modelo propuesto cumplió el objetivo planteado: superó a ambos modelos SOTA en F1-score macro y exactitud balanceada, y además redujo de forma sustancial el tiempo de entrenamiento e inferencia, lo que lo hace especialmente atractivo para su despliegue en centros de salud con recursos computacionales limitados. Entre las limitaciones del estudio destacan el uso de un único dataset público, que puede no representar la variabilidad de otros equipos de captura de imágenes, y que la poda dinámica se evaluó solo con un criterio de importancia (norma $L_1$ del gradiente), sin comparar con otros criterios de poda. Como trabajo futuro se propone validar el modelo en un segundo dataset de RD adquirido con equipos distintos, y explorar criterios de poda estructurada adicionales para determinar si es posible reducir aún más el tiempo de cómputo sin sacrificar el F1-score alcanzado.
\end{ejemplo}

\backmatter

\section*{Declaraciones}

\begin{instruccion}
Completa cada punto de forma honesta. Es muy importante que incluyas el enlace a tu repositorio de código para que el proyecto sea reproducible y verificable por tu profesor.
\end{instruccion}

\begin{ejemplo}
\begin{itemize}
\item \textbf{Disponibilidad de datos:} el conjunto APTOS 2019 es público y puede descargarse desde Kaggle (\url{https://www.kaggle.com/c/aptos2019-blindness-detection}).
\item \textbf{Disponibilidad de código:} el código completo, incluyendo scripts de entrenamiento y evaluación, está disponible en \url{https://github.com/usuario-ejemplo/rd-poda-dinamica}.
\item \textbf{Contribución de los autores:} A. Martínez diseñó la arquitectura y el algoritmo de poda; C. Rodríguez implementó los baselines y realizó los experimentos; L. Gómez realizó el análisis estadístico y la redacción del manuscrito. Todos los autores revisaron y aprobaron la versión final.
\item \textbf{Conflicto de intereses:} los autores declaran no tener conflictos de interés.
\end{itemize}
\end{ejemplo}

\section*{Referencias bibliográficas}

\begin{instruccion}
Todas las referencias (papers de los modelos SOTA, del dataset y de trabajos relacionados) deben estar en tu archivo \texttt{sn-bibliography.bib} en formato BibTeX, y citarse en el texto con \texttt{\textbackslash cite\{\}}. Ejemplo de entrada BibTeX:
\begin{verbatim}
@article{bib1,
  author  = {He, Kaiming and others},
  title   = {Deep Residual Learning for Image Recognition},
  journal = {CVPR},
  year    = {2016}
}
\end{verbatim}
En Overleaf, al usar la plantilla ``Springer Nature'', este archivo \texttt{.bib} ya viene incluido: solo agrega tus entradas ahí.
\end{instruccion}

\bibliography{sn-bibliography}

\end{document}